\section{实验}
\label{sec:exp}

\subsection{实验设置}

\paragraph{数据集。} 主数据集为 UCF-Crime，采用官方划分（$1610$ 个训练视频 / $290$ 个测试视频，含帧级异常标注）。跨域验证使用 XD-Violence（多类异常）、UBnormal（合成异常）与 MSAD（多场景多类别），三者均只做\emph{零微调外测}。为检验退化鲁棒性，我们在评测阶段对测试集注入 $4$ 类退化（低光/低对比、雨雾/大气散射、视频压缩、相机抖动与遮挡）$\times$ $3$ 个强度等级，共 $12$ 种配置，退化数据离线生成并固定随机种子。

\paragraph{评估指标。} 检测精度使用帧级 AUC 与 AP，以及时序定位 mAP@0.5（时序 IoU $\geq 0.5$ 视为命中）。解释可验证性使用四项指标：证据归因召回 $\mathrm{EAR} = |E_{pred} \cap E_{gt}| / |E_{gt}|$、反事实可反证分 $\mathrm{CFS} = 1 - \frac{1}{M}\sum_m \mathbb{1}[\hat{y}(x_m^{-}) = \text{异常}]$、幻觉率 $\mathrm{HR} = N_{unsupported}/N_{statements}$、时序一致性率 $\mathrm{TCR}$。退化鲁棒性使用相对性能保持率 $\mathrm{RPR} = M_{deg}/M_{clean}$ 与置信度校准误差 $\mathrm{ECE} = \sum_m \frac{|B_m|}{n}|acc(B_m) - conf(B_m)|$。效率指标包括可训练参数量、端到端延迟（ms / 32 帧片段）与显存占用。

\paragraph{实现细节与公平性协议。} 视觉主干为 CLIP ViT-B/16，全程冻结；DAA 注入末端 $4$ 层，$r{=}4$、$\alpha{=}8$、dropout $0.05$。优化器 AdamW，weight decay $1\times10^{-2}$；训练 $20$ 个 epoch，batch 为 $2$ 并做 $8$ 步梯度累积；AMP FP16；学习率调度为 warm-up $2$ 个 epoch 加 cosine 衰减，最小学习率 $1\times10^{-6}$。所有方法统一输入尺寸 $224\times224$、统一 $32$ 帧 / $8$~fps 片段协议、统一数据划分、统一训练轮数与统一评测脚本；阈值（门控、弃权、证据对齐）仅在验证集上确定后冻结，测试集不参与调参。结果报告三个随机种子（$42$、$3407$、$2026$）的均值与标准差。硬件为单张 RTX 4090 24\,GB、i9-13900K、64\,GB 内存，软件环境为 Ubuntu 22.04 / Python 3.10 / PyTorch 2.2 / CUDA 12.1。

\subsection{主实验结果}

表~\ref{tab:main} 给出主对比结果。\evivad{} 在检测精度上小幅领先最强对比方法（AUC $+3.2$、AP $+4.0$），但在\emph{解释可验证性}上呈现较大差异：EAR 由 $40.2$ 提升至 $68.4$，HR 由 $27.8$ 降至 $9.8$。需要强调的是，重建类、弱监督与 CLIP 零样本方法本身不输出自然语言解释，因此其 EAR / HR 记为「--」而非 $0$。图~\ref{fig:comparison} 给出同一测试样本的定性对比：基线在异常区间内分数被压低（漏报）并在非异常时段虚高（误报），而完整方法在异常区间稳定高于阈值，并输出可核验的证据引用。

\begin{figure*}[t]
\centering
\includegraphics[width=\linewidth]{comparison.png}
\caption{与 Baseline 的检测效果对比（同一测试样本，\textbf{示意}）。上排四个面板依次为输入帧、Ground Truth（红框标注「行人倒地」，GT 异常区间 $5.4$--$9.1$\,s）、Baseline 预测（漏检倒地行人并在一名正常行人处产生误报）与 \evivad{} 预测（命中并输出证据引用）。下排为同一片段的异常分数曲线：红色阴影为 GT 异常区间，虚线为判定阈值 $0.5$。图中画面与曲线均为示意，非真实实拍与真实推理结果。}
\label{fig:comparison}
\end{figure*}

\begin{table*}[t]
\centering
\small
\caption{主实验对比（UCF-Crime，结果）。``--'' 表示该方法不输出自然语言解释，故 EAR / HR 不适用。箭头 $\uparrow$ 表示越大越好，$\downarrow$ 表示越小越好。}
\label{tab:main}
\begin{tabular}{lccccccc}
\toprule
方法 & AUC $\uparrow$ & AP $\uparrow$ & mAP@0.5 $\uparrow$ & EAR $\uparrow$ & HR $\downarrow$ & 延迟 $\downarrow$ & 显存 $\downarrow$ \\
\midrule
重建类（深度自编码器族）\cite{abdalla2025vad10years}       & 74.6 & 65.1 & 18.7 & --   & --   & 46\,ms  & 1.9\,GB \\
弱监督图神经网络\cite{yun2025missiongnn}                     & 78.3 & 69.8 & 23.4 & --   & --   & 89\,ms  & 3.1\,GB \\
CLIP 潜在空间零样本\cite{zanella2024clip}                    & 75.9 & 66.5 & 20.2 & --   & --   & 74\,ms  & 2.8\,GB \\
免训练 LLM 流程\cite{zanella2024trainingfree}                & 78.9 & 70.3 & 24.6 & 40.2 & 27.8 & 131\,ms & 3.7\,GB \\
\midrule
Baseline B0（冻结 VLM + 文本打分）                            & 76.8 & 68.2 & 24.1 & 41.5 & 26.4 & 118\,ms & 3.6\,GB \\
\textbf{\evivad{}（完整方法，）}                          & \textbf{82.1} & \textbf{74.3} & \textbf{31.6} & \textbf{68.4} & \textbf{9.8} & 147\,ms & 4.1\,GB \\
\bottomrule
\end{tabular}
\end{table*}

\subsection{消融实验}

表~\ref{tab:ablation} 给出三个模块的单模块、双模块与完整组合结果。可以看到三个模块在互不重叠的维度上起作用：DAA 是精度主贡献（AUC $+2.6$，且在双模块组合中持续叠加）；EAD 是可验证性主贡献（EAR $+22.3$、CFS $+18.5$、HR $-13.8$），对精度贡献有限（$+1.1$）；DAG 在干净集上贡献最小（$+0.8$），但在退化网格上贡献最大（见下节）。三模块叠加后 AUC 达到 $82.1$，高于三个单模块增益的简单相加（$76.8+2.6+1.1+0.8 = 81.3$），说明存在弱的正向协同。

\begin{table*}[t]
\centering
\small
\caption{消融实验（UCF-Crime，结果）。$\checkmark$ 表示启用该模块，HR 越低越优。}
\label{tab:ablation}
\begin{tabular}{ccccccccccc}
\toprule
DAA & EAD & DAG & AUC $\uparrow$ & AP $\uparrow$ & mAP@0.5 $\uparrow$ & EAR $\uparrow$ & CFS $\uparrow$ & HR $\downarrow$ & TCR $\uparrow$ & 参数(M) $\downarrow$ \\
\midrule
     &     &     & 76.8 & 68.2 & 24.1 & 41.5 & 38.7 & 26.4 & 38.1 & 0.0 \\
$\checkmark$ &   &   & 79.4 & 71.1 & 26.0 & 43.2 & 40.5 & 25.1 & 39.4 & 4.7 \\
   & $\checkmark$ &   & 77.9 & 69.4 & 25.3 & 63.8 & 57.2 & 12.6 & 52.7 & 0.6 \\
   &   & $\checkmark$ & 77.6 & 69.0 & 24.9 & 43.0 & 40.2 & 24.0 & 38.9 & 0.2 \\
$\checkmark$ & $\checkmark$ &   & 80.9 & 72.6 & 28.4 & 65.7 & 58.9 & 11.2 & 54.0 & 5.3 \\
$\checkmark$ &   & $\checkmark$ & 80.3 & 72.0 & 27.5 & 44.6 & 42.1 & 22.7 & 40.2 & 4.9 \\
   & $\checkmark$ & $\checkmark$ & 79.8 & 71.4 & 27.0 & 65.1 & 58.4 & 10.9 & 53.2 & 0.8 \\
$\checkmark$ & $\checkmark$ & $\checkmark$ & \textbf{82.1} & \textbf{74.3} & \textbf{31.6} & \textbf{68.4} & \textbf{61.3} & \textbf{9.8} & \textbf{55.8} & 5.5 \\
\bottomrule
\end{tabular}
\end{table*}

\subsection{退化鲁棒性与泛化分析}

表~\ref{tab:degradation} 报告 $12$ 种退化配置下的结果。DAG 在干净集上仅带来 $+0.8$ 的 AUC，但在退化网格上平均带来 $+8.8$（$62.5 \rightarrow 71.3$）；相对性能保持率 RPR 由 $0.81$ 提升到 $0.92$，意味着退化条件下近九成性能得以保持。退化下的幻觉率同步收敛（$41.2 \rightarrow 31.5$），与「免训练方法在低光下产生幻觉语义」的既有观察一致。图~\ref{fig:curves}(a) 给出退化强度—性能曲线：基线随退化加剧急剧下滑，而 DAG 的衰减更平缓，二者差距随强度增大而扩大。图~\ref{fig:curves}(b)(c) 分别给出单模块消融的 AUC 与解释指标。

值得注意的是，DAG 属于「主要贡献体现在输入退化条件下」的模块：如果只报告干净集指标，它的价值会被严重低估。这也是我们主张把退化鲁棒性作为独立评测维度、而非附属于主指标的原因。

\begin{table}[t]
\centering
\small
\caption{退化鲁棒性（结果）。轻/中/重为三个退化强度等级，每格为该强度下 $4$ 类退化的对应取值。}
\label{tab:degradation}
\begin{tabular}{lcccc}
\toprule
退化配置 & B0 AUC $\uparrow$ & +DAG AUC $\uparrow$ & B0 HR $\downarrow$ & +DAG HR $\downarrow$ \\
\midrule
干净（无退化）      & 76.8 & 77.6 & 26.4 & 24.0 \\
低光 轻 / 中 / 重   & 66.1 / 61.3 / 55.4 & 74.2 / 70.8 / 64.6 & 34.1 / 42.7 / 51.9 & 26.8 / 31.4 / 38.6 \\
雨雾 轻 / 中 / 重   & 67.4 / 62.6 / 56.8 & 75.0 / 71.5 / 66.4 & 33.5 / 41.9 / 50.4 & 26.2 / 30.8 / 37.5 \\
压缩 轻 / 中 / 重   & 68.9 / 64.2 / 58.7 & 76.4 / 73.0 / 68.8 & 32.2 / 40.3 / 48.8 & 25.4 / 29.6 / 36.1 \\
抖动 轻 / 中 / 重   & 67.8 / 62.9 / 57.6 & 75.4 / 72.1 / 67.4 & 32.9 / 41.1 / 49.6 & 25.9 / 30.2 / 36.8 \\
\midrule
\textbf{退化网格平均} & \textbf{62.5} & \textbf{71.3} & \textbf{41.2} & \textbf{31.5} \\
RPR（相对干净集）   & 0.81 & \textbf{0.92} & -- & -- \\
\bottomrule
\end{tabular}
\end{table}

\begin{figure*}[t]
\centering
\includegraphics[width=\linewidth]{curves.png}
\caption{指标曲线（结果）。(a) 退化强度—性能曲线：横轴为干净、轻、中、重四个退化档位（每档为 $4$ 类退化的平均），基线随退化加剧急剧下滑，DAG 衰减更平缓。(b) 单模块消融的帧级 AUC。(c) 解释可验证性指标 EAR（越大越好）与 HR（越小越好）在单模块与完整方法下的对比。}
\label{fig:curves}
\end{figure*}

\subsection{效率分析与失败案例}

\evivad{} 的额外可训练参数为 $5.5$\,M（DAA $4.7$ + 证据投影头 $0.6$ + Q-Net $0.2$），端到端延迟由 $118$\,ms 升至 $147$\,ms（$+24.6\%$），显存由 $3.6$\,GB 升至 $4.1$\,GB。相对于其带来的精度与鲁棒性收益，该开销处于可接受范围；且由于视觉主干完全冻结，部署时可用「一场景一适配器」的方式替换 DAA，而不必重新分发整个模型。

我们对四类失败模式做了预登记式的边界分析：（i）\emph{门控误伤正常输入}——若 Q-Net 低估干净片段质量，$g$ 会偏向检索路径，导致干净集 AUC 下降；判据是干净集 AUC 下降不得超过 $0.3$ 个百分点，超过 $1$ 个百分点即判定该模块失败。（ii）\emph{证据对齐在人群密集片段上失效}——多目标重叠使证据槽位无法唯一对应真值，表现为密集子集 EAR 显著低于整体；我们要求报告中显式给出二者的差值。（iii）\emph{弃权阈值影响召回}——$\delta$ 过大会把正确的高置信判定一并划入弃权，因此必须报告「弃权率—召回」曲线，并把弃权样本单列而不隐藏。（iv）\emph{检索记忆缺失}——当索引为空或所有邻居相似度低于接受阈值时，融合权重退化为先验主导并输出显式的 \texttt{retrieval-miss} 标志，而非静默给出常数分数。上述四项的共同原则是：任何降级路径都必须是显式的、可观测的、可统计的。
